\scanpage{40}
\mspage{38}

par adjonction

\[
 \mathrm{Hom}_C(X,\underset{\substack{\Vert\\\psi(K')}}{K})
 =\mathrm{Hom}_{\unsure{C|C'}}(X,\psi(K'))
 \overset{\mathrm{adj.}}{\simeq}\mathrm{Hom}_{C'}(\varphi X,\deleted{\unsure{\varphi\psi}}(K'))
\]

\deleted{or $\varphi\psi(K')\xrightarrow{\sim}K'$ (adj, et le fait que $\psi$ soit pl. fid), donc} \inserted{or $\varphi(K)=\varphi\psi(K')\simeq K'$ (adj, et $\psi$ pl. fid)}

\diagram{\begin{tikzcd}[ampersand replacement=\&,column sep=small]
 \mathrm{Hom}_C(X,K)\arrow[rr,"\sim"]\arrow[dr,"{\varphi_{X,K}}"] \& \& \mathrm{Hom}_{C}(\varphi X,K')\\
 \& \mathrm{Hom}_{C'}(\varphi X,\underbrace{\varphi K}_{\substack{\deleted{\unsure{\psi\varphi K'}}\\ \Vert\,\varphi\psi(K')}})\arrow[ur,"\sim","\mathrm{adj}"'] \&
\end{tikzcd}}

\inserted{qu'on complète en}

ce qui prouve que $\varphi_{X,K}$ est iso.

(ii$''$) $\longrightarrow$ (ii$'$) \unsure{Car|car} \deleted{\unsure{la}} flèche

\diagram{\begin{tikzcd}[ampersand replacement=\&]
 \mathrm{Hom}_C(X,K)\arrow[r,"\alpha_X"]\arrow[dr,dashed] \&
 \mathrm{Hom}_{C'}(\varphi X,\varphi K)\arrow[d,dash,"\wr\ \mathrm{adj}"]\\
 \& \mathrm{Hom}_C(X,\psi\varphi K)
\end{tikzcd}}

est déduite de la flèche d'adjonction

\[
 K\longrightarrow\psi\varphi K
\]

en \unsure{prenant|posant} $\mathrm{Hom}_C(X,-)$, donc $\alpha_X$ un iso pour tout $X$ ssi cette flèche d'adjonction est iso. Donc

\[
 \mathrm{ii}\Longleftrightarrow\mathrm{ii}'\Longleftrightarrow\mathrm{ii}''
\]

D'autre part,

\[
 \mathrm{ii}''\Longrightarrow\mathrm{i}\Longrightarrow\mathrm{i}'\Longrightarrow\mathrm{i}''
\]

\margin{Car la flèche d'adjonction \deleted{$\unsure{X\to\psi\varphi X}$} $X\to\psi\varphi X$ \unsure{dans $C$ est|donc est} transformée par $\varphi$ en $\varphi X\to\underbrace{\varphi\psi}\varphi(X)$ dont le composé avec $(\varphi\psi)\varphi X\to\varphi X$}

sont tautologiques, \deleted{il reste à prouver $\mathrm{i}''\Longrightarrow\mathrm{ii}'$,} \unsure{en tenant compte|en tant} pour (i) $\Longrightarrow$ \unsure{(i)|(i$'$)} que \deleted{i.e. en supposant que pour la flèche} les flèches d'adjonction $X\xrightarrow{u}\psi\varphi X$ de $C$ sont bien dans $W$, car on a un triangle commutatif

\diagram{\begin{tikzcd}[ampersand replacement=\&]
 \varphi(X)\arrow[r,"{\varphi(u)}"]\arrow[dr,swap,"{\mathrm{id}_{\varphi(X)}}"] \&
 \varphi(\psi\varphi X)=(\varphi\psi)(\varphi X)\arrow[d,"{\wr\ \mathrm{adj.}\ \varphi\psi\to\mathrm{id}\quad\text{qui est iso car }\psi\ \text{pl. fid}}"]\\
 \& \varphi(X)
\end{tikzcd}}

donc $\varphi(u)$ bien iso. Reste à prouver $\mathrm{i}''\Longrightarrow\mathrm{ii}'$, i.e. on suppose

% archive sheet number: 39
